Srovnání \ac{GPG} podle odvětví \ac{NACE} v~ČR (\ac{ISPV}) a~longitudinálního vývoje nekorigovaného \ac{GPG} v~6 srovnávaných zemích (Eurostat) ukazuje, že odvětvový profil \ac{GPG} v~ČR je výrazně heterogenní: zatímco ve finančnictví (K) a~vyšší administrativě přesahuje \SI{25}{\percent}, ve výrobě (C) se pohybuje kolem \SIrange{10}{15}{\percent}.
Výrazná disperze \ac{GPG} napříč odvětvími při kontrole průměrného vzdělání a~pracovní doby naznačuje, že rozdíl není podmíněn pouze strukturálními faktory, ale také mírou diskrece zaměstnavatelů při individuálním sjednávání mezd --- ta je největší tam, kde chybí tarifní mzdová soustava v~\ac{KS}.
Longitudinální srovnání zemí ukazuje, že AT a~DK v~průběhu posledních 20 let \ac{GPG} snižovaly výrazněji než CZ, přičemž klíčovým strukturálním faktorem je právě přítomnost transparentních odvětvových \ac{KS} s~mzdovými tabulkami bez genderové diferenciace.
Transpozice směrnice \ac{EU} o~transparentnosti odměňování~\cite{smernice_pay_transparency_2023} v~kombinaci se zákonnou povinností publikovat mzdové tabulky ve všech registrovaných \ac{KS} jsou identifikovány jako minimálně nutné podmínky pro systémové snižování \ac{GPG} v~ČR.

Srovnání \ac{GPG} podle odvětví \ac{NACE} v~ČR (\ac{ISPV}) a~longitudinálního vývoje nekorigovaného \ac{GPG} v~6 srovnávaných zemích (Eurostat) ukazuje, že odvětvový profil \ac{GPG} v~ČR je výrazně heterogenní: zatímco ve finančnictví (K) a~vyšší administrativě přesahuje \SI{25}{\percent}, ve výrobě (C) se pohybuje kolem \SIrange{10}{15}{\percent}.
Výrazná disperze \ac{GPG} napříč odvětvími při kontrole průměrného vzdělání a~pracovní doby naznačuje, že rozdíl není podmíněn pouze strukturálními faktory, ale také mírou diskrece zaměstnavatelů při individuálním sjednávání mezd --- ta je největší tam, kde chybí tarifní mzdová soustava v~\ac{KS}.
Longitudinální srovnání zemí ukazuje, že AT a~DK v~průběhu posledních 20 let \ac{GPG} snižovaly výrazněji než CZ, přičemž klíčovým strukturálním faktorem je právě přítomnost transparentních odvětvových \ac{KS} s~mzdovými tabulkami bez genderové diferenciace.
Transpozice směrnice \ac{EU} o~transparentnosti odměňování~\cite{smernice_pay_transparency_2023} v~kombinaci se zákonnou povinností publikovat mzdové tabulky ve všech registrovaných \ac{KS} jsou identifikovány jako minimálně nutné podmínky pro systémové snižování \ac{GPG} v~ČR.

Srovnání \ac{GPG} podle odvětví \ac{NACE} v~ČR (\ac{ISPV}) a~longitudinálního vývoje nekorigovaného \ac{GPG} v~6 srovnávaných zemích (Eurostat) ukazuje, že odvětvový profil \ac{GPG} v~ČR je výrazně heterogenní: zatímco ve finančnictví (K) a~vyšší administrativě přesahuje \SI{25}{\percent}, ve výrobě (C) se pohybuje kolem \SIrange{10}{15}{\percent}.
Výrazná disperze \ac{GPG} napříč odvětvími při kontrole průměrného vzdělání a~pracovní doby naznačuje, že rozdíl není podmíněn pouze strukturálními faktory, ale také mírou diskrece zaměstnavatelů při individuálním sjednávání mezd --- ta je největší tam, kde chybí tarifní mzdová soustava v~\ac{KS}.
Longitudinální srovnání zemí ukazuje, že AT a~DK v~průběhu posledních 20 let \ac{GPG} snižovaly výrazněji než CZ, přičemž klíčovým strukturálním faktorem je právě přítomnost transparentních odvětvových \ac{KS} s~mzdovými tabulkami bez genderové diferenciace.
Transpozice směrnice \ac{EU} o~transparentnosti odměňování~\cite{smernice_pay_transparency_2023} v~kombinaci se zákonnou povinností publikovat mzdové tabulky ve všech registrovaných \ac{KS} jsou identifikovány jako minimálně nutné podmínky pro systémové snižování \ac{GPG} v~ČR.

Srovnání \ac{GPG} podle odvětví \ac{NACE} v~ČR (\ac{ISPV}) a~longitudinálního vývoje nekorigovaného \ac{GPG} v~6 srovnávaných zemích (Eurostat) ukazuje, že odvětvový profil \ac{GPG} v~ČR je výrazně heterogenní: zatímco ve finančnictví (K) a~vyšší administrativě přesahuje \SI{25}{\percent}, ve výrobě (C) se pohybuje kolem \SIrange{10}{15}{\percent}.
Výrazná disperze \ac{GPG} napříč odvětvími při kontrole průměrného vzdělání a~pracovní doby naznačuje, že rozdíl není podmíněn pouze strukturálními faktory, ale také mírou diskrece zaměstnavatelů při individuálním sjednávání mezd --- ta je největší tam, kde chybí tarifní mzdová soustava v~\ac{KS}.
Longitudinální srovnání zemí ukazuje, že AT a~DK v~průběhu posledních 20 let \ac{GPG} snižovaly výrazněji než CZ, přičemž klíčovým strukturálním faktorem je právě přítomnost transparentních odvětvových \ac{KS} s~mzdovými tabulkami bez genderové diferenciace.
Transpozice směrnice \ac{EU} o~transparentnosti odměňování~\cite{smernice_pay_transparency_2023} v~kombinaci se zákonnou povinností publikovat mzdové tabulky ve všech registrovaných \ac{KS} jsou identifikovány jako minimálně nutné podmínky pro systémové snižování \ac{GPG} v~ČR.

\input{texparts/python/rscp_gender_gap}



